% Options for packages loaded elsewhere
\PassOptionsToPackage{unicode}{hyperref}
\PassOptionsToPackage{hyphens}{url}
\PassOptionsToPackage{dvipsnames,svgnames,x11names}{xcolor}
%
\documentclass[
  11pt,
]{article}
\usepackage{amsmath,amssymb}
\usepackage{iftex}
\ifPDFTeX
  \usepackage[T1]{fontenc}
  \usepackage[utf8]{inputenc}
  \usepackage{textcomp} % provide euro and other symbols
\else % if luatex or xetex
  \usepackage{unicode-math} % this also loads fontspec
  \defaultfontfeatures{Scale=MatchLowercase}
  \defaultfontfeatures[\rmfamily]{Ligatures=TeX,Scale=1}
\fi
\usepackage{lmodern}
\ifPDFTeX\else
  % xetex/luatex font selection
\fi
% Use upquote if available, for straight quotes in verbatim environments
\IfFileExists{upquote.sty}{\usepackage{upquote}}{}
\IfFileExists{microtype.sty}{% use microtype if available
  \usepackage[]{microtype}
  \UseMicrotypeSet[protrusion]{basicmath} % disable protrusion for tt fonts
}{}
\makeatletter
\@ifundefined{KOMAClassName}{% if non-KOMA class
  \IfFileExists{parskip.sty}{%
    \usepackage{parskip}
  }{% else
    \setlength{\parindent}{0pt}
    \setlength{\parskip}{6pt plus 2pt minus 1pt}}
}{% if KOMA class
  \KOMAoptions{parskip=half}}
\makeatother
\usepackage{xcolor}
\usepackage[margin=25mm]{geometry}
\usepackage{listings}
\newcommand{\passthrough}[1]{#1}
\lstset{defaultdialect=[5.3]Lua}
\lstset{defaultdialect=[x86masm]Assembler}
\usepackage{longtable,booktabs,array}
\usepackage{calc} % for calculating minipage widths
% Correct order of tables after \paragraph or \subparagraph
\usepackage{etoolbox}
\makeatletter
\patchcmd\longtable{\par}{\if@noskipsec\mbox{}\fi\par}{}{}
\makeatother
% Allow footnotes in longtable head/foot
\IfFileExists{footnotehyper.sty}{\usepackage{footnotehyper}}{\usepackage{footnote}}
\makesavenoteenv{longtable}
\setlength{\emergencystretch}{3em} % prevent overfull lines
\providecommand{\tightlist}{%
  \setlength{\itemsep}{0pt}\setlength{\parskip}{0pt}}
\setcounter{secnumdepth}{-\maxdimen} % remove section numbering
\emergencystretch=3em
\AtBeginDocument{\Urlmuskip=0mu plus 2mu}
\lstset{breaklines=true,basicstyle=\small\ttfamily,columns=fullflexible,keepspaces=true}
\setlength{\tabcolsep}{4pt}
\usepackage{iftex}
\ifXeTeX
  \usepackage{newunicodechar}
  \makeatletter
  \newcommand{\DeclareUnicodeCharacter}[2]{%
    \begingroup
    \uccode`\~="#1\relax
    \uppercase{\def\temp{~}}%
    \expandafter\endgroup
    \expandafter\newunicodechar\expandafter{\temp}{#2}%
  }
  \makeatother
\fi
\ifdefined\DeclareUnicodeCharacter
\DeclareUnicodeCharacter{2192}{\ensuremath{\to}}
\DeclareUnicodeCharacter{2190}{\ensuremath{\leftarrow}}
\DeclareUnicodeCharacter{2194}{\ensuremath{\leftrightarrow}}
\DeclareUnicodeCharacter{21D2}{\ensuremath{\Rightarrow}}
\DeclareUnicodeCharacter{21A6}{\ensuremath{\mapsto}}
\DeclareUnicodeCharacter{2212}{\ensuremath{-}}
\DeclareUnicodeCharacter{2264}{\ensuremath{\leq}}
\DeclareUnicodeCharacter{2265}{\ensuremath{\geq}}
\DeclareUnicodeCharacter{2260}{\ensuremath{\neq}}
\DeclareUnicodeCharacter{2113}{\ensuremath{\ell}}
\DeclareUnicodeCharacter{03A3}{\ensuremath{\Sigma}}
\DeclareUnicodeCharacter{03A6}{\ensuremath{\Phi}}
\DeclareUnicodeCharacter{0393}{\ensuremath{\Gamma}}
\DeclareUnicodeCharacter{221E}{\ensuremath{\infty}}
\DeclareUnicodeCharacter{2208}{\ensuremath{\in}}
\DeclareUnicodeCharacter{2209}{\ensuremath{\notin}}
\DeclareUnicodeCharacter{211A}{\ensuremath{\mathbb{Q}}}
\DeclareUnicodeCharacter{211D}{\ensuremath{\mathbb{R}}}
\DeclareUnicodeCharacter{2124}{\ensuremath{\mathbb{Z}}}
\DeclareUnicodeCharacter{2115}{\ensuremath{\mathbb{N}}}
\DeclareUnicodeCharacter{2200}{\ensuremath{\forall}}
\DeclareUnicodeCharacter{2203}{\ensuremath{\exists}}
\DeclareUnicodeCharacter{2286}{\ensuremath{\subseteq}}
\DeclareUnicodeCharacter{00D7}{\ensuremath{\times}}
\DeclareUnicodeCharacter{03B4}{\ensuremath{\delta}}
\DeclareUnicodeCharacter{2202}{\ensuremath{\partial}}
\DeclareUnicodeCharacter{03B2}{\ensuremath{\beta}}
\DeclareUnicodeCharacter{03B1}{\ensuremath{\alpha}}
\DeclareUnicodeCharacter{03C6}{\ensuremath{\phi}}
\DeclareUnicodeCharacter{03B3}{\ensuremath{\gamma}}
\DeclareUnicodeCharacter{00B9}{\ensuremath{^1}}
\DeclareUnicodeCharacter{00B2}{\ensuremath{^2}}
\DeclareUnicodeCharacter{2070}{\ensuremath{^0}}
\DeclareUnicodeCharacter{2080}{\ensuremath{_0}}
\DeclareUnicodeCharacter{2081}{\ensuremath{_1}}
\DeclareUnicodeCharacter{2082}{\ensuremath{_2}}
\fi
\ifLuaTeX
  \usepackage{selnolig}  % disable illegal ligatures
\fi
\IfFileExists{bookmark.sty}{\usepackage{bookmark}}{\usepackage{hyperref}}
\IfFileExists{xurl.sty}{\usepackage{xurl}}{} % add URL line breaks if available
\urlstyle{same}
\hypersetup{
  pdftitle={PFP VIII: Conservative Grammar Growth and Source-Bound Evidence in the MSS Pipeline},
  pdfauthor={Jeremy H. Carroll},
  colorlinks=true,
  linkcolor={Maroon},
  filecolor={Maroon},
  citecolor={Blue},
  urlcolor={Blue},
  pdfcreator={LaTeX via pandoc}}

\title{PFP VIII: Conservative Grammar Growth and Source-Bound Evidence
in the MSS Pipeline}
\author{Jeremy H. Carroll}
\date{2026-10-02 · Version 0.1.0 · Research methods paper ·
Author-review revision}

\begin{document}
\makeatletter
\begin{center}
{\LARGE \@title\par}
\vspace{0.8em}
{\normalsize \@author\par}
\vspace{0.5em}
{\small \@date\par}
\end{center}
\makeatother

\textbf{Code baseline:}
\nolinkurl{52cb5457f6c485ba548cfd292f6012a941ed22cf}

\hypertarget{abstract}{%
\subsection{Abstract}\label{abstract}}

Adding a new symbol, checking a finite example, and proving a semantic
theorem are different achievements. This paper studies how the Proposal
Fidelity Protocol (PFP, historically the Symbol Game) records those
distinctions in the MSS grammar and its evidence pipeline. MSS-N denotes
stage N in the packet-language catalog. The implemented context maps
preserve the single old packet at the MSS-0 to MSS-1 transition and
return a typed non-descent receipt for new structure. That limited round
trip establishes injectivity on its declared domain, but does not
establish a general tower or a conservative extension of a theory. We
connect this example to computed stage grades, finite Lean declarations,
and a versioned resolver that separates historical Git bytes from
installed resources. A worked flow fixture shows why matching objective
values can coexist with an invalid certificate. We then apply the same
standard to a challenge corpus and to recorded model runs. A census at a
stated commit sorts 50 challenges into delivered solver packages, model
measurements, and specifications, and corrects earlier counts. In
recorded runs of two small models on one graph, the correct repair
appeared only when the prompt contained the repair values, and the
host's acceptance rule could also be met by discarding the claims. The
contribution is a source-pinned methods case study: explicit
preservation obligations, adversarial controls, and evidence-custody
checks that a reader can reproduce without accepting a stage name or a
success label as proof. The deposit includes the actual initial parser
and maps, two retained working reports, and all eight K4 traces with
full revision and hash bindings. A 16-question study companion turns
these examples into predictions and changed-input checks.

\hypertarget{what-must-survive-an-extension}{%
\subsection{1. What must survive an
extension?}\label{what-must-survive-an-extension}}

Suppose a language initially has one legal packet, written as
\passthrough{\lstinline!<>!}. A later language adds named packets such
as \passthrough{\lstinline!<0>!} and \passthrough{\lstinline!<1>!}. Two
requirements are easy to confuse. First, the old packet should remain
identifiable: storing it in the new language and reading it back should
recover the same object. Second, its old meaning and all claims made
about it should remain justified. The first requirement alone does not
establish the second.

This distinction matters when a research system grows by adding stages.
A stage number may organize a catalog, but it cannot certify that the
stage's advertised mathematics has been implemented. Likewise, a
successful Lean compilation establishes propositions expressed in Lean;
it does not automatically establish the broader meaning suggested by a
file's title.

We examine four layers in the pinned source: actual packet syntax and
context maps, stage-specific executable checks, formal declarations and
their build evidence, and retrieval of the exact bytes cited by a
result. The central question is operational: what may a reader conclude
at each layer, and what observation would invalidate that conclusion?

This is a methods and audit paper. It does not introduce an infinite
mathematical tower, infer physical laws from a grammar, or report a
controlled study of model reasoning.

\hypertarget{the-packet-language-and-the-implemented-round-trip}{%
\subsection{2. The packet language and the implemented round
trip}\label{the-packet-language-and-the-implemented-round-trip}}

The
\href{../../008_pfp_mssn_tower_challenges_and_retraction_benchmarks/docs/source_archive/mss0_mss1_52cb5457/grammar.py}{archived
grammar module} uses the following recursive syntax:

\begin{lstlisting}
packet ::= diff | atom | app
diff   ::= <>
atom   ::= <word>
app    ::= op<packet*>
word   ::= one or more ASCII characters from a-z or 0-9
\end{lstlisting}

Packets are self-delimiting. An application can contain adjacent child
packets without commas. The empty distinction
\passthrough{\lstinline!<>!} has its own Diff constructor; it is not an
Atom with an empty name. Parsing produces a structured value; a separate
level-admission predicate decides which packets are active in a
particular context.

For the context-map module, the only MSS-0 packet is Diff. The module
\href{../../008_pfp_mssn_tower_challenges_and_retraction_benchmarks/docs/source_archive/mss0_mss1_52cb5457/updown_maps.py}{updown\_maps.py}
implements precisely u\_0 and d\_1. The upward map first checks that its
input is active at level zero, then returns that same packet. The
downward map returns an old packet unchanged and returns a
NonDescentReceipt containing the original packet otherwise.

Let \(S_0=\{\mathrm{Diff}\}\), and let \(S_1\) be the intended active
MSS-1 packet domain. Its interface is better represented by

\[
u_0:S_0\longrightarrow S_1,\qquad
d_1:S_1\longrightarrow S_0\sqcup\mathrm{Refusal},
\qquad d_1(u_0(x))=x.
\]

The disjoint alternative matters: a refusal is not an ordinary old
packet. Calling the down-map a total function from \(S_1\) to \(S_0\)
would erase the distinction that the implementation deliberately
preserves.

\begin{longtable}[]{@{}
  >{\raggedright\arraybackslash}p{(\columnwidth - 4\tabcolsep) * \real{0.3333}}
  >{\raggedright\arraybackslash}p{(\columnwidth - 4\tabcolsep) * \real{0.3333}}
  >{\raggedright\arraybackslash}p{(\columnwidth - 4\tabcolsep) * \real{0.3333}}@{}}
\toprule\noalign{}
\begin{minipage}[b]{\linewidth}\raggedright
Input operation
\end{minipage} & \begin{minipage}[b]{\linewidth}\raggedright
Result
\end{minipage} & \begin{minipage}[b]{\linewidth}\raggedright
What it establishes
\end{minipage} \\
\midrule\noalign{}
\endhead
\bottomrule\noalign{}
\endlastfoot
\passthrough{\lstinline!d\_1(u\_0(parse("<>")))!} & Diff & Old packet is
recovered exactly \\
\passthrough{\lstinline!u\_0(parse("<0>"))!} & ValueError & Non-old
input violates the upward precondition \\
\passthrough{\lstinline!d\_1(parse("<1>"))!} & NonDescentReceipt
retaining \passthrough{\lstinline!<1>!} & New structure is not silently
converted into Diff \\
\end{longtable}

These are the concrete cases in the
\href{https://github.com/dbccpage/the_symbol_game/blob/52cb5457f6c485ba548cfd292f6012a941ed22cf/research/mssn_sn_tower/code/python/tests/test_mssn_updown_and_cells.py}{context-map
tests}. The d\_1 function itself accepts a typed Packet and tests
old-level membership; it does not independently certify that every
supplied packet belongs to \(S_1\). A caller needing that guarantee must
run the level validator.

\textbf{Left-inverse fact.} If \(d(u(x))=x\) on an old domain, then
\(u\) is injective there: applying \(d\) to \(u(x)=u(y)\) gives \(x=y\).
In the implemented transition the old domain has one element, and the
identity follows directly from the two function branches. This is a
useful API invariant, with a deliberately small mathematical scope.

The first part of the package-local test file
\href{../../008_pfp_mssn_tower_challenges_and_retraction_benchmarks/code/tests/test_008_mssn_tower_retraction.py}{\nolinkurl{test_008_mssn_tower_retraction.py}}
is a separate teaching toy. It changes a record's level field and uses a
string-prefix rule to return None for a new term. Its example does not
implement the domain maps above, and None is not their typed refusal. We
retain that toy as historical material; the paper's implementation
claims refer to the actual modules.

\hypertarget{why-injectivity-is-weaker-than-semantic-conservativity}{%
\subsection{3. Why injectivity is weaker than semantic
conservativity}\label{why-injectivity-is-weaker-than-semantic-conservativity}}

A conservative extension of a theory must preserve its old consequences:
if the extended theory proves a sentence in the old language, that
sentence must already follow from the original theory. A round trip on
packet objects does not mention derivability, truth, or an
interpretation. Verified compilation shows the stronger kind of
statement: the correctness theorem of the CompCert back-end relates the
observable behaviors of a source program and its compiled code (Leroy,
2009). A syntactic round trip makes no such comparison.

For a simple counterexample, keep every old term unchanged while adding
a new inference rule that declares an unproved old-language proposition
\(P\) to be an axiom. The identity inclusion of terms still has a left
inverse. The extended proof system nevertheless has a new old-language
consequence. Consequently, injective syntax cannot justify semantic
conservativity.

An extension that intends to preserve meaning needs additional data: an
interpretation for both languages, a compatibility statement on embedded
old terms, and a theorem that the allowed new rules do not create
unsupported old conclusions. A stronger operational claim might also
require preservation of accepted transitions, failure states, and
observable outputs. None of these follows by an increase in an integer
stage field alone.

The audited code baseline does not supply a family of u\_n and d\_n maps
with these properties for all stages. Its source explicitly restricts
the concrete context-map implementation to MSS-0 and MSS-1. Larger
catalogs and proposed obligations remain useful design material, but
should be cited as such.

One namespace fact does generalize. The
\href{https://github.com/dbccpage/the_symbol_game/blob/52cb5457f6c485ba548cfd292f6012a941ed22cf/research/mssn_sn_tower/code/python/registry.py}{registry}
allocates eight consecutive proposal opcodes starting at

\[
b(N)=128+8(N-22),\qquad N\geq22.
\]

For \(N<M\), \(b(N)+7<b(M)\); hence these integer intervals are
disjoint. The source opcodes are below 128, so the proposal blocks also
avoid the declared source namespace. Python's unbounded integers support
this formula; a fixed-width wire encoding would require an additional
upper bound. The loaded catalog covers stages zero through 100, not
every natural number. Neither the arithmetic formula nor the catalog
constructs those stages' semantics.

\hypertarget{evidence-grades-come-from-checks}{%
\subsection{4. Evidence grades come from
checks}\label{evidence-grades-come-from-checks}}

The
\href{https://github.com/dbccpage/the_symbol_game/blob/52cb5457f6c485ba548cfd292f6012a941ed22cf/research/mssn_sn_tower/code/python/stage_status.py}{stage-status
module} uses the grades NONE, SYNTAX, and COMPUTED. It does not grant a
grade because a level lies below a chosen threshold.

A SYNTAX check round-trips the stage's own packet through serialization,
admits it at its declared stage, and refuses it one stage early. For
stage zero, the negative probe is an interior glyph that should not yet
be admitted. These checks are evidence of the selected syntax boundary.
They do not prove the entire parser correct for every input.

A COMPUTED binding identifies a checker and exact fixtures. The checker
must accept its positive example and refuse its stage-specific negative
example and empty input; floating-point evidence is excluded. Registered
target checkers additionally bind checker bytes and the stage-row
digest, and exercise near misses, mutations of each deciding clause, and
applicable probes with their expected refusal codes.

For example, the finite-section binding covers \(\{a,b\}\) and
\(\{b,c\}\). Assignments \((a=0,b=1)\) and \((b=1,c=0)\) glue to
\((a=0,b=1,c=0)\). Replacing the second value at \(b\) by zero creates a
conflict and must fail. This validates a finite dictionary operation; it
does not establish a general sheaf theorem.

Registering a checker says which executable is responsible for a stage.
Running it supplies the current grade. Tests that pass on a small
fixture family still leave untested inputs, specification mistakes, and
common errors shared by a checker and its fixtures. The grade is
therefore a bounded evidence label, not a universal soundness theorem.

This status interface never assigns LEAN or GATED. Lean evidence belongs
to a separate pipeline. A stale comment that no symbol-game axiom report
has been captured must not be read as a repository-wide inventory: the
Level Engine has its own report artifacts and custody requirements.

\hypertarget{a-formal-fixture-that-exposes-an-invalid-certificate}{%
\subsection{5. A formal fixture that exposes an invalid
certificate}\label{a-formal-fixture-that-exposes-an-invalid-certificate}}

At the code baseline, the Level Engine directory contains 52 payload
modules plus its axiom-report module. Their audited source is in
\href{https://github.com/dbccpage/the_symbol_game/blob/52cb5457f6c485ba548cfd292f6012a941ed22cf/lean4/research_programs/level_engine/README.md}{\nolinkurl{lean4/research_programs/level_engine}}.
Their statements must be read individually.

Consider
\href{https://github.com/dbccpage/the_symbol_game/blob/52cb5457f6c485ba548cfd292f6012a941ed22cf/lean4/research_programs/level_engine/FlowCertificate.lean}{\mbox{\nolinkurl{FlowCertificate.lean}}}.
It defines a three-edge directed cycle, unit penalties, edge loads
\(c=(1,0,0)\), and zero vertex levels. Its primal absolute-residual cost
is one. With flows \(y=(1,1,1)\), the dual dot product is also one,
capacity bounds hold, and the flow balances at every vertex.

Now change only the flow to \(y=(1,0,0)\). The dot product remains
\(c^\mathsf{T}y=1\), so the reported gap is still zero. Capacity bounds
also remain true. But the divergence, defined as inflow minus outflow,
is \((-1,1,0)\). This vector cannot serve as a feasible dual
circulation.

The declarations explicitly prove that the own-negative fixture has
zeroGap equal to true and balanced equal to false. Thus a theorem about
one arithmetic clause can be true even when the proposed certificate
fails as a whole.
\href{https://github.com/dbccpage/the_symbol_game/blob/23893fc44162555fa12ceda3e1665982e3d7a14e/atlas/coda_atlas/symbol_game/001_pfp_discrete_cochain_retraction_and_toy_involutions/paper.md}{Paper
001} supplies the general weak-duality argument explaining why every
certificate clause is needed.

The file uses exact-rational definitions and kernel-decided finite
propositions. It does not by itself prove a general optimization theorem
over all graphs. Similarly, the following source readings place limits
on several suggestive filenames. Each item gives what the file states or
constructs and the further conclusion it does not establish.

\begin{itemize}
\tightlist
\item
  \href{https://github.com/dbccpage/the_symbol_game/blob/52cb5457f6c485ba548cfd292f6012a941ed22cf/lean4/Level16IndObject.lean}{\mbox{\nolinkurl{Level16IndObject.lean}}}:
  natural-number capacities \(2^{n+1}\) increase and exceed every finite
  bound. Not established: a category, a colimit construction, or a
  theorem about embeddings of that colimit.
\item
  \href{https://github.com/dbccpage/the_symbol_game/blob/52cb5457f6c485ba548cfd292f6012a941ed22cf/lean4/Level21DiscreteRegge.lean}{\mbox{\nolinkurl{Level21DiscreteRegge.lean}}}:
  concrete rotation commutator has zero scalar exponent sum and
  nonidentity matrix product. Not established: general discrete Regge
  dynamics.
\item
  \href{https://github.com/dbccpage/the_symbol_game/blob/52cb5457f6c485ba548cfd292f6012a941ed22cf/lean4/Level22FibredDependentContexts.lean}{\mbox{\nolinkurl{Level22FibredDependentContexts.lean}}}:
  reindexing copies a record's fields; identity and composition hold for
  that operation. Not established: a Grothendieck fibration with its
  universal lifting property.
\item
  \href{https://github.com/dbccpage/the_symbol_game/blob/52cb5457f6c485ba548cfd292f6012a941ed22cf/lean4/Level23SheafDescent.lean}{\mbox{\nolinkurl{Level23SheafDescent.lean}}}:
  restrictions of a global function are compatible; conflicting
  restrictions prevent amalgamation. Not established: a general
  existence-and-uniqueness gluing theorem.
\end{itemize}

The revised Level21 file expressly withdraws itself as evidence for its
earlier MSS-021 claim and formalizes a finite counterexample instead.
That correction must survive later summaries. A build result cannot
restore the withdrawn interpretation.

Copies of these five files are vendored in this package's
\passthrough{\lstinline!lean4/!} folder
(\href{../../008_pfp_mssn_tower_challenges_and_retraction_benchmarks/lean4/lakefile.lean}{Lake
file}), together with the ExactRat library that
\mbox{\nolinkurl{FlowCertificate.lean}} imports. Each copy is
byte-identical to its source at the code baseline apart from a two-line
provenance comment. The retained author receipts record their
compilation with Lean 4.34.0, and the
\href{../../008_pfp_mssn_tower_challenges_and_retraction_benchmarks/lean4/axiom_report.txt}{axiom
report} lists the axioms of all 39 theorems the five files declare.
Thirty-eight depend on no axiom beyond propext and Quot.sound. The
theorem stage\_capacity\_monotonic in the Level16 file also depends on
Classical.choice. The two-axiom policy of Section 6 belongs to the Level
Engine audit and governs only its payload modules; the Level16 file is
not one of them.

The broader local
\href{../../008_pfp_mssn_tower_challenges_and_retraction_benchmarks/lean4/full_axiom_report.txt}{report}
lists 576 environment declarations from ten modules, including ExactRat
and generated or private declarations. This is not a count of manuscript
theorems. The retained
\href{../../008_pfp_mssn_tower_challenges_and_retraction_benchmarks/docs/release_validation_2026_10_02.json}{October
2 author receipt} records a successful local build and this broader
census. Review checked its 13 protected source hashes, the full census
source and report, and the census tool against the deposited bytes; all
match. A fresh Lean replay in the present review host is blocked: the
official, digest-verified Lean 4.34.0 binary exits before startup with
``failed to locate application''. The retained source bindings remain
inspectable, but this review does not claim an independent kernel pass.

\hypertarget{fresh-formal-evidence-needs-custody}{%
\subsection{6. Fresh formal evidence needs
custody}\label{fresh-formal-evidence-needs-custody}}

A proof assistant checks a proposition under its definitions and
assumptions. Lean's documentation explains the role of extensionality,
quotients, and choice in that logical environment.
\href{https://lean-lang.org/theorem_proving_in_lean4/Axioms-and-Computation/}{Official
Lean reference}.

The Level Engine
\href{https://github.com/dbccpage/the_symbol_game/blob/52cb5457f6c485ba548cfd292f6012a941ed22cf/operations/use_cases/level_engine/lean_build.py}{build
driver} pins Lean 4.34.0 and records source hashes, return codes, and
whether the sources remained stable. Its
\href{https://github.com/dbccpage/the_symbol_game/blob/52cb5457f6c485ba548cfd292f6012a941ed22cf/operations/use_cases/level_engine/lean_audit.py}{axiom-audit
driver} requires a successful build receipt matching the current source
bytes and toolchain before running a fresh report.

The report must cover the complete generator-declared theorem census
without missing, extra, or duplicate names. The local policy permits
only propext and Quot.sound. It rejects dependencies such as sorryAx,
Classical.choice, or a custom unproved axiom, and compares the fresh
report with the canonical report. This is the package's chosen policy; a
theorem using classical choice is not thereby a mathematical falsehood.

At the pinned baseline, the driver checks source hashes before and after
Lean returns. Review found a later publication gap: a source edit during
census inspection can be recorded under an earlier successful audit
because the publisher recomputes its expected hashes. A bounded mocked
reproduction returned success and published the changed source; the
downstream snapshot reader correctly rejected its mismatched build/audit
identities.

The local correction in the
\href{https://github.com/dbccpage/the_symbol_game/blob/23893fc44162555fa12ceda3e1665982e3d7a14e/atlas/coda_atlas/symbol_game/verification/source_delta.patch}{verification
patch} carries the audited source map through the receipt and snapshot
writer, refuses a post-audit edit, and validates copied bytes against
that map. The regression tests exercise this control logic with fake
Lean responses. They do not establish a fresh Lean build, filesystem
locking, or protection against an edit that is made and reverted between
observations.

The
\href{https://github.com/dbccpage/the_symbol_game/blob/52cb5457f6c485ba548cfd292f6012a941ed22cf/operations/tests/level_engine/test_lean.py}{custody
tests} use isolated fake tool responses to check these failure paths.
They include a changed source after a success receipt, an absent
receipt, an incomplete declaration census, and disallowed axioms even
when report replacement is requested. They test the control logic. They
are not fresh Lean kernel elaborations of the 52 payloads.

Neither this review nor the retained local build rebuilt the 52 external
Level Engine payloads. The local 39-theorem selection and the
576-declaration environment census are different scopes; neither
certifies those payloads. A new full Level Engine result requires its
own source-bound execution receipt.

The deposit's optional named Lean gate also compares the complete
expected declaration names and axiom sets with the replay. An empty
successful process, a missing, extra or duplicate declaration, or a
changed dependency is refused. Mocked process tests validate this gate
logic and do not compile a theorem.

\hypertarget{resolve-the-cited-bytes-not-a-convenient-replacement}{%
\subsection{7. Resolve the cited bytes, not a convenient
replacement}\label{resolve-the-cited-bytes-not-a-convenient-replacement}}

An old evidence path may no longer exist in the working directory.
Substituting a current file with the same name can quietly change the
evidence while leaving the citation looking intact.

The
\href{https://github.com/dbccpage/the_symbol_game/blob/52cb5457f6c485ba548cfd292f6012a941ed22cf/symbol_game/infrastructure/tower_evidence/level_engine/resolver.py}{versioned
resolver} has two deliberately separate read modes:

\begin{longtable}[]{@{}
  >{\raggedright\arraybackslash}p{(\columnwidth - 4\tabcolsep) * \real{0.3333}}
  >{\raggedright\arraybackslash}p{(\columnwidth - 4\tabcolsep) * \real{0.3333}}
  >{\raggedright\arraybackslash}p{(\columnwidth - 4\tabcolsep) * \real{0.3333}}@{}}
\toprule\noalign{}
\begin{minipage}[b]{\linewidth}\raggedright
Mode
\end{minipage} & \begin{minipage}[b]{\linewidth}\raggedright
Required source
\end{minipage} & \begin{minipage}[b]{\linewidth}\raggedright
Missing-data behavior
\end{minipage} \\
\midrule\noalign{}
\endhead
\bottomrule\noalign{}
\endlastfoot
v1 & Explicit repository, full 40-character revision, canonical relative
path & \nolinkurl{HISTORICAL_BYTES_UNAVAILABLE} \\
v2 & Installed package resource under the active v2 namespace &
\nolinkurl{ACTIVE_RESOURCE_UNAVAILABLE} \\
\end{longtable}

The historical mode invokes Git's object reader for the exact
revision-and-path blob; it does not fall back to the working tree. The
active mode uses package resources and refuses a historical revision
argument. Both reject absolute paths, parent traversal, backslashes,
empty segments, and other noncanonical forms. An optional SHA-256 check
rejects a mismatch. A caller that omits this digest has not obtained
that additional content check.

The
\href{https://github.com/dbccpage/the_symbol_game/blob/52cb5457f6c485ba548cfd292f6012a941ed22cf/operations/tests/level_engine/test_evidence_resolver.py}{resolver
tests} create a temporary Git repository, commit bytes containing line
endings and a zero byte, remove the working file, and read the original
bytes through v1. They also check wrong digests, unavailable revisions,
missing repository context, noncanonical paths, and misuse of the v2
revision argument. Git's documented object-content operation supplies
the underlying retrieval primitive.
\href{https://git-scm.com/docs/git-cat-file}{Git cat-file
documentation}.

This mechanism preserves byte identity under its stated assumptions. It
does not establish who authored the bytes, whether a checker is correct,
or whether a mathematical claim follows from them. The caller must still
bind each cited artifact to the proposition and source revision being
reviewed.

\hypertarget{a-benchmark-label-is-not-an-observation}{%
\subsection{8. A benchmark label is not an
observation}\label{a-benchmark-label-is-not-an-observation}}

The historical
\href{https://github.com/dbccpage/the_symbol_game/blob/52cb5457f6c485ba548cfd292f6012a941ed22cf/projects/mssn_sn_symgam/challenges/run_150_tests_suite.py}{150-test
runner} contains a concrete evidence failure: some absent-receipt cases
receive VERIFIED, caught provider exceptions can also receive VERIFIED,
and some route or metric values are fixed records rather than
measurements from the named experiment. The aggregate counts PASSED and
VERIFIED alike.

Therefore an aggregate 150/150 record from that procedure cannot
establish 150 independently executed successful scientific checks. Nor
can it support the former manuscript's claims of ``perfect governed
resolution'', ``zero ungoverned completion'', or a measured retraction
frequency. Those claims are withdrawn. The recorded spend and duration
are not adopted as validated benchmark results.

A sound replacement requires recording which operation actually ran,
preserving raw outputs, distinguishing missing and failed execution from
success, and applying an independent outcome oracle to each experimental
claim. The controlled-evaluation design belongs to
\href{https://github.com/dbccpage/the_symbol_game/blob/23893fc44162555fa12ceda3e1665982e3d7a14e/atlas/coda_atlas/symbol_game/003_pfp_host_admitted_failure_detectors/paper.md}{Paper
003}. Sections 9 and 10 apply the same standard to the challenge corpus
and to recorded model runs.

\hypertarget{the-challenge-corpus-at-a-stated-commit}{%
\subsection{9. The challenge corpus at a stated
commit}\label{the-challenge-corpus-at-a-stated-commit}}

At the census revision,
\nolinkurl{projects/mssn_sn_symgam/challenges/} holds
challenge specifications written against one protocol,
\href{https://github.com/dbccpage/the_symbol_game/blob/a623568eb0f853b968bc419bf06dfc6af256ab02/projects/mssn_sn_symgam/challenges/01_challenge_protocol.md}{SG-CP
v1}. This paper audits C001 to C050, in five subfolders: mechanism (C001
to C010), cohomology (C011 to C020), quantum information (C021 to C030),
physics (C031 to C040), and formal mathematics (C041 to C050). The
folder also holds C051 to C100 and G01 to G25, which are outside this
paper.

The protocol asks each challenge for a participant packet and a judge
packet. The judge packet holds hidden cases, relabel and gauge controls,
poisoned but schema-valid cases, an exact checker, a second checker
written separately, and a mutation suite: corrupted outputs that the
checker must reject. The protocol admits only open-weight models with at
most 70B total parameters. It sets a per-arm spending ceiling by tier,
from \$0.10 at tier T1 to \$5.00 at tier T5, enforced before each call.
It asks for a committed pre-registration and exact McNemar tests on
paired cases. It requires exact arithmetic and allows floating-point
checks only as labeled secondary evidence.

\hypertarget{census}{%
\subsubsection{9.1 Census}\label{census}}

We classify every challenge from the tree of commit
\nolinkurl{a623568eb0f853b968bc419bf06dfc6af256ab02} (2026-09-28). The
challenge folder is read at that historical revision; it is not
identified with today's main branch. The three categories are disjoint:

\begin{itemize}
\tightlist
\item
  \textbf{A. Delivered solver with two checkers}: a reference solver,
  two checker modules in the judge packet, and executable checks.
\item
  \textbf{B. Model measurement deliverable}: the specification calls for
  a measurement of a language model rather than a reference solver.
  Membership does not assert that a measurement was completed.
\item
  \textbf{C. Specification only}: a specification exists and no
  reference solver is delivered.
\end{itemize}

\begin{longtable}[]{@{}llll@{}}
\toprule\noalign{}
Subfolder & A & B & C \\
\midrule\noalign{}
\endhead
\bottomrule\noalign{}
\endlastfoot
Mechanism & none & C001, C005 & C002 to C004, C006 to C010 \\
Cohomology & C011 to C020 & none & none \\
Quantum information & C021 to C030 & none & none \\
Physics & C031 to C040 & none & none \\
Formal mathematics & C041 to C050 & none & none \\
\end{longtable}

Category A therefore holds 40 of the 50 challenges. C051, outside the
audited range, is a forty-first delivered package. The
\href{../../008_pfp_mssn_tower_challenges_and_retraction_benchmarks/docs/challenge_census.md}{census
file} lists every Category A folder with its two checker modules, its
number of test files, hidden cases, and mutation files, and its retained
scoring receipt. File counts are not pytest collection counts. C047 has
one test-named file, \nolinkurl{judge_packet/test_cases_runner.py},
which prints scores for its cases using both checkers, without
assertions or a failing exit status; it does not check the mutation
files. C048 to C051 have no pytest test file in the inventory. Their
scoring programs must be read for their actual scope; their presence
does not establish that every hidden case and mutation was executed. The
C014 receipt records that its scoring was invalidated and a rescore is
pending. No scoring receipt is retained for C015 to C020 or C045 to
C051. Membership in Category A records what a folder contains; it does
not certify that a package is correct.

The previous revision listed C001 and C005 in both Categories A and B,
and it listed C012 to C020, C022 to C030, C032 to C040, and C043 to C050
as specifications without solvers. That list was wrong even at the code
baseline: packages for C012 to C020 and C032 to C040 were already
present at commit \passthrough{\lstinline!52cb5457!}. Packages for C022
to C030 and C041 to C051 were added after it, on 2026-09-27 and
2026-09-28.

\hypertarget{retained-author-checks}{%
\subsubsection{9.2 Retained author
checks}\label{retained-author-checks}}

The September 28 author ledger records five test-file runs from three
Category A packages, one file per command, with Python 3.12.10 and
pytest 8.3.5. All five are recorded as exit status 0. The present review
inspected the historical census and retained source; it did not rerun
these remote challenge packages.

\begin{longtable}[]{@{}
  >{\raggedright\arraybackslash}p{(\columnwidth - 2\tabcolsep) * \real{0.5000}}
  >{\raggedright\arraybackslash}p{(\columnwidth - 2\tabcolsep) * \real{0.5000}}@{}}
\toprule\noalign{}
\begin{minipage}[b]{\linewidth}\raggedright
Package and test file
\end{minipage} & \begin{minipage}[b]{\linewidth}\raggedright
Result
\end{minipage} \\
\midrule\noalign{}
\endhead
\bottomrule\noalign{}
\endlastfoot
C011 \nolinkurl{judge_packet/mutation_suite/test_checker_mutations.py} &
41 passed: an inventory check, and 40 mutation files each rejected by
both checkers \\
C011 \nolinkurl{tests/test_dual_checkers.py} & 5 passed \\
C021 \nolinkurl{judge_packet/mutation_suite/test_checker_mutations.py} &
41 passed: an inventory check and 40 mutation files \\
C031 \nolinkurl{judge_packet/mutation_suite/test_checker_mutations.py} &
47 passed: an inventory check, 45 mutation files, and a check on refusal
mutations \\
C031 \nolinkurl{tests/test_reference_solver.py} & 9 passed \\
\end{longtable}

These runs and the retained files correct earlier figures. C011 has 46
judge cases (32 base cases, 8 invariance controls, and 6 poisoned cases)
and 40 mutation files, not 32. Its checkers verify an exact dual
certificate for minimum weighted \(L_1\) coboundary repair: the primal
cost \(\sum_e w_e\,|c_e-(\delta_0\phi)_e|\) equals \(\sum_e c_e z_e\)
for a circulation \(z\) with zero divergence at every vertex and
\(|z_e|\le w_e\). C031 has 45 mutation files, not 36, and the earlier
count of 71 tests is not traceable to the retained files; it is
withdrawn. Its reference-solver tests check that no chiral solution
exists for \(n=3\) or \(n=4\) within the charge bounds 3, 5, 10, and 12,
and that for \(n=5\) with \(|q_i|\le 9\) the enumerator returns exactly
one canonical solution, \((-9,-5,-1,7,8)\). Up to sign and order this is
\((1,5,-7,-8,9)\), for which \(1+5-7-8+9=0\) and
\(1+125-343-512+729=0\). The earlier figure of 38 end-to-end tests for
C021 was the number of test functions in the shared folder
\passthrough{\lstinline!challenges/e2e\_tests/!}, which this revision
did not run; it is withdrawn. C021's second checker includes a
floating-point eigensolver that the package labels as secondary
evidence.

We did not run the C041 checkers, since each case needs a Lean call. By
source inspection, they refuse a fixed list of tokens that includes
\passthrough{\lstinline!sorry!}, \passthrough{\lstinline!axiom!},
\passthrough{\lstinline!unsafe!}, \passthrough{\lstinline!set\_option!},
and \passthrough{\lstinline!native\_decide!}, and they accept a proof
only when its axioms lie in \{propext, Quot.sound\}. Proofs by
\passthrough{\lstinline!native\_decide!} depend on the additional axiom
\passthrough{\lstinline!Lean.ofReduceBool!}. The package records Lean
4.34.1 as its toolchain; the earlier text named 4.34.0.

\hypertarget{category-b-model-measurements}{%
\subsubsection{9.3 Category B: model
measurements}\label{category-b-model-measurements}}

C001 contains a harness that runs Qwen2.5-0.5B-Instruct with real
weights, and a 200-case McNemar module that describes itself as a
synthetic fixture with injected baseline failures. No measurement result
from the harness is retained in the repository: it writes its report
only to an untracked run folder, and it lists the full 200-case sweep as
its next test. The earlier statement that C001 evaluated distractor-key
masking across 200 tasks is withdrawn.

C005 compares prompt and output formats on 50 graph-consistency cases
(25 consistent, 25 with a planted cycle defect) with one model,
\nolinkurl{nousresearch/hermes-3-llama-3.1-405b}, a fine-tune of Llama
3.1 405B, at temperature 0. The host's exact integer checker scored
every answer: the verdict had to match the case label, and a CONSISTENT
verdict needed a potential that satisfies every edge. The results come
from the retained
\href{../../008_pfp_mssn_tower_challenges_and_retraction_benchmarks/docs/source_archive/c005_working_report.md}{archived
C005 working report}, which describes two September 27 runs and retains
its September 29 correction notice. Its old aggregate wording is not
adopted here.

\begin{longtable}[]{@{}
  >{\raggedright\arraybackslash}p{(\columnwidth - 6\tabcolsep) * \real{0.2500}}
  >{\raggedright\arraybackslash}p{(\columnwidth - 6\tabcolsep) * \real{0.2500}}
  >{\raggedright\arraybackslash}p{(\columnwidth - 6\tabcolsep) * \real{0.2500}}
  >{\raggedright\arraybackslash}p{(\columnwidth - 6\tabcolsep) * \real{0.2500}}@{}}
\toprule\noalign{}
\begin{minipage}[b]{\linewidth}\raggedright
Run
\end{minipage} & \begin{minipage}[b]{\linewidth}\raggedright
Arm
\end{minipage} & \begin{minipage}[b]{\linewidth}\raggedright
Prompt and output format
\end{minipage} & \begin{minipage}[b]{\linewidth}\raggedright
Correct
\end{minipage} \\
\midrule\noalign{}
\endhead
\bottomrule\noalign{}
\endlastfoot
Three-arm run & 1 & Relations in words; verdict first & 25/50 \\
Three-arm run & 2 & Relations in words; numeric scratchpad before the
verdict & 36/50 \\
Three-arm run & 3 & Precomputed spanning-tree potentials and chords;
chord residuals only & 19/50 \\
Hybrid run & 2 & As in the three-arm run & 38/50 \\
Hybrid run & 4 & Factorized input, a canary field, and a scratchpad; the
canary field must also be right & 4/50 \\
\end{longtable}

In the paired comparisons below, \(b\) counts cases correct only under
the first arm and \(c\) counts cases correct only under the second. We
computed the exact two-sided binomial McNemar p-value from \(b\) and
\(c\). For \(n=b+c>0\), we use

\[
p=\min\left(1,\;2\,2^{-n}\sum_{k=0}^{\min(b,c)}\binom{n}{k}\right),
\]

and set \(p=1\) when \(n=0\). These are descriptive paired comparisons,
without a multiple-comparison adjustment.

\begin{longtable}[]{@{}
  >{\raggedright\arraybackslash}p{(\columnwidth - 8\tabcolsep) * \real{0.2000}}
  >{\raggedright\arraybackslash}p{(\columnwidth - 8\tabcolsep) * \real{0.2000}}
  >{\raggedright\arraybackslash}p{(\columnwidth - 8\tabcolsep) * \real{0.2000}}
  >{\raggedright\arraybackslash}p{(\columnwidth - 8\tabcolsep) * \real{0.2000}}
  >{\raggedright\arraybackslash}p{(\columnwidth - 8\tabcolsep) * \real{0.2000}}@{}}
\toprule\noalign{}
\begin{minipage}[b]{\linewidth}\raggedright
Comparison and run
\end{minipage} & \begin{minipage}[b]{\linewidth}\raggedright
b
\end{minipage} & \begin{minipage}[b]{\linewidth}\raggedright
c
\end{minipage} & \begin{minipage}[b]{\linewidth}\raggedright
Exact p
\end{minipage} & \begin{minipage}[b]{\linewidth}\raggedright
Approximate p
\end{minipage} \\
\midrule\noalign{}
\endhead
\bottomrule\noalign{}
\endlastfoot
Arm 1 vs Arm 2, three-arm run & 6 & 17 & \(\frac{145499}{4194304}\) &
0.0347 \\
Arm 1 vs Arm 3, three-arm run & 6 & 0 & \(\frac{1}{32}\) & 0.0313 \\
Arm 2 vs Arm 3, three-arm run & 20 & 3 & \(\frac{1}{2048}\) & 0.00049 \\
Arm 2 vs Arm 4, hybrid run & 34 & 0 & \(2^{-33}\) &
\(1.2\times10^{-10}\) \\
\end{longtable}

Within each run the counts and the paired tables agree: 25 − 6 + 17 =
36, 25 − 6 = 19, 36 − 20 + 3 = 19, and 38 − 34 = 4. The previous
revision reported Arm 2 as 37/50 (74.0\%), the midpoint of the two runs,
and set it beside tables taken from both runs; no run produced 37/50.
The per-case rows were written to an untracked run folder and are not in
the repository, so the counts above are transcribed from the report and
could not be recomputed from per-case data here.

The C005 runs also depart from the protocol. The
\href{https://github.com/dbccpage/the_symbol_game/blob/a623568eb0f853b968bc419bf06dfc6af256ab02/projects/mssn_sn_symgam/challenges/a_mechanism/c005_hosted_decode_controls_bakeoff.md}{C005
specification} and SG-CP v1 both limit models to open weights with at
most 70B total parameters and set a ceiling of \$0.10 per arm. The runs
used a 405B model and a ceiling of \$2.00 per benchmark run. They are
reported here as a deviation from the protocol, not as a conforming C005
result. On this one model and case set, the two factorized-input arms
scored below the verdict-first arm; that observation does not support a
general statement about prompt or schema complexity.

\hypertarget{recorded-retraction-traces}{%
\subsection{10. Recorded retraction
traces}\label{recorded-retraction-traces}}

\hypertarget{eight-k4-traces}{%
\subsubsection{10.1 Eight K4 traces}\label{eight-k4-traces}}

The eight retained traces are archived in
\href{../../008_pfp_mssn_tower_challenges_and_retraction_benchmarks/docs/source_archive/README.md}{\nolinkurl{docs/source_archive/k4_traces/}}.
The accompanying
\href{../../008_pfp_mssn_tower_challenges_and_retraction_benchmarks/docs/source_archive/k4_benchmark_report.md}{working
benchmark report} (last changed 2026-09-25) records a deterministic
replay through the host. Its custody table gives SHA-256 digests for
eight traces, one for each pair of model and condition. We recomputed
all eight digests and they match. There is one trace per cell, so every
rate below is 0/1 or 1/1. The earlier figure of 28 trials, and the
per-cell counts of 4 and 3 derived from it, are not traceable to
retained files and are withdrawn.

All eight runs use one fixture: the complete graph \(K_4\) on vertices
a, b, c, d, with cycle rank 3, OpenRouter as the provider, temperature
0, and 12 turns. The model asserts the six edge claims and then proposes
repair moves. After each move the host recomputes its defect value
\(\Phi_1\), which is zero exactly when the current claims admit a
potential. A repair move replaces the whole claim set with the claims it
lists, an omitted edge counting as zero, and the host admits it unless
\(\Phi_1\) increases. A session is accepted when the model verifies at
\(\Phi_1=0\) or when the session ends at \(\Phi_1=0\). The correct
repair keeps the three tree edges and sets the chords to c→a = −3, b→d =
−3, and c→d = −5. The four conditions differ in what the host's prompt
supplies once all six claims are asserted:

\begin{longtable}[]{@{}
  >{\raggedright\arraybackslash}p{(\columnwidth - 2\tabcolsep) * \real{0.5000}}
  >{\raggedright\arraybackslash}p{(\columnwidth - 2\tabcolsep) * \real{0.5000}}@{}}
\toprule\noalign{}
\begin{minipage}[b]{\linewidth}\raggedright
Condition
\end{minipage} & \begin{minipage}[b]{\linewidth}\raggedright
What the prompt supplies beyond the graph
\end{minipage} \\
\midrule\noalign{}
\endhead
\bottomrule\noalign{}
\endlastfoot
Baseline & The sum around each of the three basis cycles, and the
one-word guidance ``qualitative'' \\
Partially masked & The spanning-tree potentials and a formula for each
chord value; the model must do three subtractions \\
Scratchpad & As partially masked, plus a worked example whose arithmetic
yields exactly −3, −3, and −5 \\
Hinted & The target value of each chord and a complete recommended
repair move \\
\end{longtable}

The scratchpad and hinted prompts therefore contain the repair values.
They are assisted conditions: a pass shows that a model copied or
reproduced supplied values, not that it found the repair unaided.

\begin{longtable}[]{@{}llllll@{}}
\toprule\noalign{}
Model & Condition & Turns & Loop count & Final \(\Phi_1\) & Outcome \\
\midrule\noalign{}
\endhead
\bottomrule\noalign{}
\endlastfoot
Qwen2.5-7B-Instruct & Baseline & 12 & 6 & 8 & Not accepted \\
Qwen2.5-7B-Instruct & Partially masked & 12 & 5 & 6 & Not accepted \\
Qwen2.5-7B-Instruct & Scratchpad & 9 & 1 & 0 & Accepted \\
Qwen2.5-7B-Instruct & Hinted & 9 & 1 & 0 & Accepted \\
Llama-3.1-8B-Instruct & Baseline & 12 & 5 & 5 & Not accepted \\
Llama-3.1-8B-Instruct & Partially masked & 12 & 4 & 7 & Not accepted \\
Llama-3.1-8B-Instruct & Scratchpad & 12 & 4 & 0 & Accepted \\
Llama-3.1-8B-Instruct & Hinted & 12 & 4 & 0 & Accepted \\
\end{longtable}

The loop count follows the counting rule of the benchmark script, which
the deposit does not include; its retained last version is
\nolinkurl{research/mssn_sn_tower/demos/demo_1/benchmark_open_weights_nullspace.py}
at
\href{https://github.com/dbccpage/the_symbol_game/blob/7b0d2903a13ceada014f572e4a3f9073ad37f95a/research/mssn_sn_tower/demos/demo_1/benchmark_open_weights_nullspace.py}{commit
\nolinkurl{7b0d2903a13ceada014f572e4a3f9073ad37f95a}}. A turn counts
when the host classifies a move as stationary or refuses it for raising
the defect, or when the model asserts an edge after turn 6. The rule
therefore also counts Qwen's seventh assertion, which was needed because
it asserted the first edge twice. The count is a label from that rule,
not a count of identical repeated proposals. Each trace also records the
host's replies to every turn except the last, and those replies are
consistent with the retained replay table. The final row has no stored
post-move host result. The present review verified all eight trace
hashes and inspected their proposals; the terminal defect values above
remain the historical report's replay, not a fresh replay or a final
field read from each raw trace.

Three details of the traces qualify the table. First, both accepted Qwen
runs proposed the correct repair at turn 8 and verified at turn 9.
Second, both Llama runs in the assisted conditions proposed the correct
repair at turn 7, moved away from it, and proposed it again later (turn
11 with the scratchpad, turn 12 with the hint); both were accepted by
the end-of-session rule. The scratchpad run's last move set the three
chords to zero and omitted the tree edges; the host admits such a move,
which leaves \(\Phi_1=0\) on a claim set of zeros. Third, Llama's turn-9
move in the baseline condition was of the same kind. The host labeled it
an exact repair, although it replaced the claims with zeros rather than
repairing them; later assertions restored the defect.

These records support one narrow statement. On this fixture, neither
model proposed the correct chord values when the prompt left the three
subtractions to it, and both models proposed them when the prompt
contained them. Acceptance was decided by the host's exact check of
\(\Phi_1\), not by the model. That check scores the claims that a move
proposes, not their agreement with the claims asserted earlier, so an
acceptance label is weaker than a repair of the asserted claims.

The previous revision said that hinted prompts reduced loops to at most
1.0 and that hinted runs converged within an average of 8.8 to 11.0
turns. With one trace per cell and Llama's hinted loop count of 4, those
figures are withdrawn. We also do not describe the hinted condition as
supplying only a boundary vector: it supplied the full repair. The same
report contains further traces for Qwen2.5-72B-Instruct and
Llama-3.3-70B-Instruct on the complete graphs \(K_4\), \(K_5\), and
\(K_6\); this paper does not analyze them.

\hypertarget{hamiltonian-traces}{%
\subsubsection{10.2 Hamiltonian traces}\label{hamiltonian-traces}}

The previous revision reported that the engine reached
exact-cancellation terminals on three two- and four-site Hamiltonians
that it called stoquastic, and that it correctly emitted a
sign-obstruction terminal for a three-site triangular Hubbard plaquette.
The retained traces are in
\nolinkurl{research/mssn_sn_tower/demos/demo_2/traces/} and in Paper
001's \passthrough{\lstinline!results/!} folder. The repository's own
\href{https://github.com/dbccpage/the_symbol_game/blob/52cb5457f6c485ba548cfd292f6012a941ed22cf/research/mssn_sn_tower/demos/anti_morphism_observations.md}{audit
of those runs} and the
\href{https://github.com/dbccpage/the_symbol_game/blob/52cb5457f6c485ba548cfd292f6012a941ed22cf/research/mssn_sn_tower/demos/demo_2/2_demo.md}{Demo
2 notes} found that the historical runners set the defect to zero after
one asserted matrix entry and accepted a named symmetry without checking
the proposed map. Their exact-cancellation terminals are therefore false
acceptances, not certificates, and the claim is withdrawn.

The triangular trace,
\nolinkurl{demo2_r5_frustrated_triangular_1790208154.jsonl}, records a
different pattern. The host refused the model's sign-free claim with a
message stating that the triangle is an odd cycle, told the model to
declare an obstruction, and recorded the obstruction terminal that the
model then declared. The fact behind that refusal is that a triangle has
no bipartition: no assignment of signs \(\eta\) to the three sites
satisfies \(\eta_u\eta_v=-1\) on every edge. That is a statement about
one family of site gauges. The prompt declared hopping and on-site terms
but no particle number, basis ordering, or path representation, so the
trace supports no classification of the Hamiltonian's sign problem. The
source of the rule that the host applied is not retained with the trace.
The baseline's check for a related question,
\href{https://github.com/dbccpage/the_symbol_game/blob/52cb5457f6c485ba548cfd292f6012a941ed22cf/symbol_game/infrastructure/adapters/mssn_research/sign_switching.py}{sign\_switching.py},
takes a real symmetric matrix in a declared basis and decides whether
some diagonal ±1 change of basis makes every off-diagonal entry
nonpositive; when none exists, it returns a cycle whose positive entries
have odd parity. The earlier sentence about diamond graphs and
``infinite 2-cycles'' had no retained receipt at the stated scope and is
withdrawn.

\hypertarget{the-package-runner}{%
\subsubsection{10.3 The package runner}\label{the-package-runner}}

The package's
\href{../../008_pfp_mssn_tower_challenges_and_retraction_benchmarks/code/retraction_benchmark_runner.py}{runner}
is a smaller, single-shot tool, separate from the recorded sessions
above. It asks whether a potential exists on six fixed graphs under four
conditions. Each host verdict has a certificate: a potential satisfying
every edge, or a closed walk with nonzero signed claim sum. A potential
telescopes to zero on every closed walk, so the latter certificate rules
out all potentials. The verifier also accepts walks loaded back from
JSON and refuses malformed or altered walks.

Scoring accepts an exact CONSISTENT or CONTRADICTION verdict; malformed
explicit labels, conflicting labels, and substring matches yield
UNKNOWN. Before querying the proposer, the runner refuses a fixture
whose declared label disagrees with the host. Each version-2 execution
receipt binds the complete input graph, exact prompt, raw reply,
provider label, measurements, host certificate, and declared
runner-source digest. Its public SHA-256 detects changes relative to a
trusted expected digest; it does not authenticate a provider or prove a
certificate, since an editor can recompute it.

The hinted prompt gives the host's edge residuals, which decide the
answer. The offline stub reads the prompt; it matches 3/6, 6/6, 3/6 and
6/6 under baseline, hinted, scratchpad and exact-oracle conditions.
These are plumbing controls, not model measurements. The local
scratchpad asks for generic subtraction reasoning; it is different from
the historical K4 scratchpad's answer-bearing example. A
repeated-proposal helper counts the longest consecutive run of repeated
comparisons, not a total pooled across separate runs. It is not the
historical K4 loop rule.

Offline mode opens no network connection or credential file. Outputs go
to an explicit directory or a new temporary directory outside the
package. Live mode requires an environment key and reports no run here.
Its nominal \$0.10 budget uses a fixed \$0.001 next-call estimate and
reported cost; it is not a guaranteed bound on provider charges.

\hypertarget{reproduction-and-practical-use}{%
\subsection{11. Reproduction and practical
use}\label{reproduction-and-practical-use}}

Use Python 3.12 from this package directory, one test file per command.
Install the versions in the review receipt if reproducibility of these
checks matters:

\begin{lstlisting}
py -3.12 -m pip install pytest==9.1.1 ruff==0.14.13 mypy==1.19.1
py -3.12 -m pytest code/tests/test_008_mssn_tower_retraction.py -q
py -3.12 -m pytest code/tests/test_008_deposit_gates.py -q
py -3.12 -m pytest code/tests/test_008_source_bound_examples.py -q
py -3.12 -m ruff check code
py -3.12 -m mypy
py -3.12 code/retraction_benchmark_runner.py --mock --condition hinted --output-dir <new-directory>
py -3.12 tools/replay_lean_axiom_census.py --check .
py -3.12 tools/replay_lean_axiom_census.py --replay .
\end{lstlisting}

On Unix systems, use the same commands with the Python 3.12 interpreter
in place of \passthrough{\lstinline!py -3.12!}. Lean replay requires a
working Lean 4.34.0 toolchain; a failed startup is a blocked replay, not
a mathematical result. The source example tests import the archived
initial grammar and maps in an isolated namespace. They also recompute
the flow fixture and the four exact McNemar values, and compare
certificates for all 125 integer triangle assignments with edge values
from −2 to 2 against an independent bounded potential search. This is a
finite test domain, not a proof of a general graph algorithm.

The
\href{../../008_pfp_mssn_tower_challenges_and_retraction_benchmarks/_lev_examiner_theoretical_minimum_quiz/quiz.md}{quiz}
and
\href{../../008_pfp_mssn_tower_challenges_and_retraction_benchmarks/_lev_examiner_theoretical_minimum_quiz/study_guide.md}{worked
study guide} connect each calculation to its source, scope limit and
changed-input control.

The old repository-root commands refer to the code baseline, whose
directory layout has since changed. To replay those historical checks,
use a separate checkout at that full commit; use the census commit for
the challenge packages. Historical Level Engine generation, build and
audit likewise need that checkout, its dependencies and vendor inputs.
They are separate from the standalone deposit checks and are not
executed by this review.

The approach separates declaring a syntax extension, checking a finite
instance, and asserting general semantic preservation.

\hypertarget{scope-of-the-claims}{%
\subsection{12. Scope of the claims}\label{scope-of-the-claims}}

This paper supports the following statements.

\begin{enumerate}
\def\labelenumi{\arabic{enumi}.}
\tightlist
\item
  The MSS-0 to MSS-1 context maps form a left inverse on a one-element
  domain and return a typed non-descent receipt for new packets. A left
  inverse gives injectivity; it does not give a conservative extension.
\item
  At commit \passthrough{\lstinline!a623568e!}, 40 of the challenges
  C001 to C050 are delivered solver packages with two checker modules, 2
  call for model measurements, and 8 are specifications only. The
  retained author runs for C011, C021 and C031 record rejection of their
  40, 40 and 45 mutation files respectively.
\item
  The C005 report's counts, read per run, agree with its paired tables,
  and the exact p-values follow from its discordant counts. The runs
  departed from the protocol's model-size and budget limits.
\item
  In eight recorded K4 traces, one per model and condition, the two 7B
  and 8B models proposed the correct repair exactly in the conditions
  whose prompts contained the repair values, and the host's acceptance
  rule could also be met by a move that replaced the claims with zeros.
\item
  Retained source-bound author receipts record a Lean 4.34.0 build of
  the local formal package. Its 39 selected theorems and the broader
  576-declaration census have different scopes. This review verifies
  those source bindings and reports its fresh local Lean replay as
  blocked.
\end{enumerate}

This paper does not support the following statements.

\begin{enumerate}
\def\labelenumi{\arabic{enumi}.}
\tightlist
\item
  That an infinite tower, or a conservative extension for every stage,
  exists.
\item
  That any Category A package is correct beyond its retained receipt and
  the historical author runs, or that the Category C specifications are
  implemented.
\item
  Any measured C001 result, any general claim about prompt formats from
  C005, or any reading of the C005 runs as a protocol-conforming result.
\item
  Unassisted repair by the models, or any rate estimated from single
  traces.
\item
  Any classification of the sign problem for the Hamiltonian traces, or
  any polynomial-time solution of general sign problems.
\end{enumerate}

The narrowest defensible claim is this. In the recorded runs, the host's
exact check of \(\Phi_1\) decided which proposals were accepted; the
prompts changed only what the models proposed. On one K4 fixture, the
two small models proposed the correct repair only when the prompt
contained the repair values.

The next test named by the previous revision, reference solvers for C012
and C022, is no longer open: both packages are delivered and hold
retained scoring receipts. The open next test is an unassisted
comparison: several seeds on the K4 fixture with prompts that give the
potentials but not the chord values, against a matched control, under a
pre-registered analysis, together with publication of per-case C005 rows
so that its tables can be recomputed.

\hypertarget{limitations}{%
\subsection{13. Limitations}\label{limitations}}

The census records folder contents at one historical commit. Section 9.2
retains the author's five challenge-test runs; this review did not
replay them. The C005 counts rest on a report whose per-case rows are
not retained. The K4 traces are single runs on one fixture at
temperature 0, and their loop count comes from one script's rule. The
Hamiltonian traces support no physical conclusion. The Level Engine
payloads were not rebuilt. Retained local Lean source identities match,
but the fresh review-host replay is blocked. The 16-question curriculum
is a self-study design, not a measured educational intervention.

This paper assigns no authorization and changes no jurisdiction. It
carries no operational authority or institutional standing.

\hypertarget{verification-for-this-revision}{%
\subsection{Verification for this
revision}\label{verification-for-this-revision}}

The package verifier (\passthrough{\lstinline!verify\_deposit.py!})
checks the manifest and the two-way file census, and
\passthrough{\lstinline!tools/rebuild\_release.py --check!} checks the
release archive against the manifests. The
\href{../../008_pfp_mssn_tower_challenges_and_retraction_benchmarks/docs/publication_review_2026_10_02.md}{publication
review} and
\href{../../008_pfp_mssn_tower_challenges_and_retraction_benchmarks/docs/review_validation_2026_10_02.json}{source-bound
review receipt} record this review's offline tests, lint, typing, source
comparisons, PDF build and blocked fresh Lean replay. They distinguish
executed checks from retained historical outcomes. The older receipts
keep their original dates and source hashes; unchanged Lean bindings
remain relevant, while their old runner and teaching-material hashes do
not describe the revised package.

\hypertarget{references}{%
\subsection{References}\label{references}}

\begin{itemize}
\tightlist
\item
  Avigad, Jeremy, Leonardo de Moura, Soonho Kong, and Sebastian Ullrich,
  with contributions from the Lean Community. ``Axioms and
  Computation.'' Chapter 12 of \emph{Theorem Proving in Lean 4}.
  Official online book; accessed 2026-09-28.
  \href{https://lean-lang.org/theorem_proving_in_lean4/Axioms-and-Computation/}{Chapter};
  \href{https://lean-lang.org/theorem_proving_in_lean4/}{book and
  authors}.
\item
  Git project. ``git-cat-file --- Provide contents or details of
  repository objects.'' \emph{Git Reference Manual}, online edition
  (release 2.56.0 when accessed on 2026-09-28).
  \href{https://git-scm.com/docs/git-cat-file}{Manual page}.
\item
  Leroy, Xavier (2009). ``A Formally Verified Compiler Back-end.''
  \emph{Journal of Automated Reasoning}, 43(4), pp.~363--446.
  \href{https://doi.org/10.1007/s10817-009-9155-4}{\nolinkurl{doi:10.1007/s10817-009-9155-4}};
  \href{https://arxiv.org/abs/0902.2137}{author preprint,
  arXiv:0902.2137}.
\end{itemize}

\end{document}
